
\documentclass[10pt]{article}
\usepackage[utf8]{inputenc}
\usepackage[T1]{fontenc}
\usepackage{amsmath, amssymb, xcolor, geometry, enumitem, multicol, tcolorbox}

\definecolor{mainblue}{RGB}{0, 102, 204}
\geometry{a4paper, margin=1.2cm, top=1cm, bottom=3cm, footskip=1.2cm}

\begin{document}
\enlargethispage{2.5cm}

% Modern Header
\begin{tcolorbox}[colback=mainblue!5, colframe=mainblue, sharp corners, boxrule=0.5pt]
    \small \textbf{Unit:} undefined \hfill \textbf{Name:} \rule{3cm}{0.4pt} \\
    \small \textbf{Topic:} Algebra 1: Polynomials \hfill \textbf{Date:} \rule{2cm}{0.4pt} \\
    \centering \textbf{\large Level 2: Foundation}
\end{tcolorbox}

\vspace{0.2cm}

% MCQ Section
\begin{multicols}{2}
\begin{enumerate}[label=\textbf{Q\arabic*.}, leftmargin=*, itemsep=6pt, parsep=0pt]

  \item \small What is the degree of the polynomial $3x^2 + 2x - 4$?
  \begin{enumerate}[label=(\Alph*), leftmargin=0.6cm, nosep, itemsep=2pt]
    \item 1
    \item 2
    \item 3
    \item 4
    \item 5
  \end{enumerate}

  \item \small Which of the following is a polynomial? 
  \begin{enumerate}[label=(\Alph*), leftmargin=0.6cm, nosep, itemsep=2pt]
    \item $x^2 + 3x = 10$
    \item $x^2 + 3x + 1$
    \item $\frac{1}{x} + 2x^2$
    \item $x^3 - 2x^2 + x + 1$
    \item $\sqrt{x} + 2x^2$
  \end{enumerate}

  \item \small What is the coefficient of $x$ in the polynomial $2x^2 + 5x - 3$?
  \begin{enumerate}[label=(\Alph*), leftmargin=0.6cm, nosep, itemsep=2pt]
    \item 2
    \item 5
    \item 3
    \item 1
    \item 0
  \end{enumerate}

  \item \small Simplify the expression: $2x^2 + 3x + 4x^2$
  \begin{enumerate}[label=(\Alph*), leftmargin=0.6cm, nosep, itemsep=2pt]
    \item $6x^2 + 3x$
    \item $5x^2 + 3x$
    \item $2x^2 + 7x$
    \item $6x^2 + 7x$
    \item $5x^2 + 5x$
  \end{enumerate}

  \item \small Factor out the greatest common factor from $6x^2 + 12x$
  \begin{enumerate}[label=(\Alph*), leftmargin=0.6cm, nosep, itemsep=2pt]
    \item $6x(x+2)$
    \item $6x^2(2x+1)$
    \item $6x^2(x+2)$
    \item $6x(x+1)$
    \item $3(2x^2 + 4x)$
  \end{enumerate}

  \item \small What is the value of $x$ that makes the polynomial $x^2 + 4x + 4$ equal to zero?
  \begin{enumerate}[label=(\Alph*), leftmargin=0.6cm, nosep, itemsep=2pt]
    \item -4
    \item -3
    \item -2
    \item 0
    \item -1
  \end{enumerate}

  \item \small Which one of the following polynomials has a degree of 3?
  \begin{enumerate}[label=(\Alph*), leftmargin=0.6cm, nosep, itemsep=2pt]
    \item $2x^2 + x - 3$
    \item $3x^3 - 2x^2 + x - 1$
    \item $x^2 + 3x - 4$
    \item $5x - 2$
    \item $x^4 + 2x^2 - 1$
  \end{enumerate}

  \item \small What is the result of adding $2x^2 + 3x - 4$ and $x^2 - 2x - 1$?
  \begin{enumerate}[label=(\Alph*), leftmargin=0.6cm, nosep, itemsep=2pt]
    \item $3x^2 + x - 3$
    \item $3x^2 + x - 5$
    \item $x^2 + x - 5$
    \item $3x^2 + 5x - 3$
    \item $x^2 - x - 3$
  \end{enumerate}

\end{enumerate}
\end{multicols}

\vspace{-0.2cm}
\hrule
\vspace{0.2cm}
\textbf{\small Open Ended Questions (FRQ)}
\begin{enumerate}[label=\textbf{Q\arabic*.}, start=9, itemsep=5pt, topsep=0pt]

  \item \small Draw a diagram to show how the polynomial $x^2 + 2x + 1$ can be factored as $(x + 1)^2$. Explain why this is a special factoring case. \vspace{2cm}

  \item \small Explain why $2x^2 + 5x - 3$ is a polynomial, but $\frac{1}{x} + 2x^2$ is not. Include definitions and examples to support your explanation. \vspace{2cm}

  \item \small Simplify the expression: $(2x^2 + 3x) + (x^2 - 2x)$. Draw a step-by-step solution and explain the reasoning behind combining like terms. \vspace{2cm}

  \item \small Factor out the greatest common factor from the expression $12x^3 + 18x^2$. Show your work and explain why this is an important step in simplifying polynomials. \vspace{2cm}

\end{enumerate}

\vfill

% Chef's Tips Footer
\begin{tcolorbox}[colback=blue!5, colframe=mainblue!50, title=\small \textbf{Chef's Tips}, fonttitle=\bfseries, bottom=1mm]
\begin{itemize}[noitemsep, topsep=2pt, leftmargin=0.5cm]
  \item \scriptsize Emphasize the importance of understanding the definition of a polynomial and how to identify the degree, terms, and coefficients.\item \scriptsize Practice combining like terms and factoring out the greatest common factor to simplify polynomials.\item \scriptsize Use visual aids and diagrams to help students understand the relationship between polynomials and their factored forms.
\end{itemize}
\end{tcolorbox}

\end{document}
  